\section{Fourier Transformation}

\subsection{Definition and Inverse}

\subsubsection{Continuous-time Fourier transform}
\[
X(\omega)=\int_{-\infty}^{\infty}x(t)e^{-j\omega t}\,dt.
\]
The Fourier transform maps a continuous-time signal to angular frequency in rad/s.

\subsubsection{Inverse Fourier transform}
\[
x(t)=\frac{1}{2\pi}\int_{-\infty}^{\infty}X(\omega)e^{j\omega t}\,d\omega.
\]
The inverse transform reconstructs the time signal from the complete frequency spectrum.

\subsubsection{Magnitude and phase spectrum}
\[
X(\omega)=|X(\omega)|e^{j\angle X(\omega)}.
\]
Magnitude describes frequency content amplitude. Phase describes time alignment of the frequency components.

\subsection{Fourier Transform Pairs}

\subsubsection{Impulse and constant}
\[
\delta(t)\ \longleftrightarrow\ 1,
\qquad
1\ \longleftrightarrow\ 2\pi\delta(\omega).
\]
An impulse contains equal weight at all frequencies. A constant signal has all spectral mass at zero frequency.

\subsubsection{Complex exponential and sinusoid}
\[
e^{j\omega_0t}\ \longleftrightarrow\ 2\pi\delta(\omega-\omega_0),
\]
\[
\cos\omega_0t\ \longleftrightarrow\ \pi[\delta(\omega-\omega_0)+\delta(\omega+\omega_0)].
\]
Periodic sinusoids transform into impulses at their positive and negative angular frequencies.

\subsubsection{Causal exponential}
\[
e^{-at}u(t)\ \longleftrightarrow\ \frac{1}{a+j\omega},\qquad a>0.
\]
The one-sided decaying exponential produces a first-order frequency response.

\subsubsection{Rectangular pulse and sinc spectrum}
\[
x(t)=
\begin{cases}
1, & |t|<T/2,\\
0, & |t|>T/2,
\end{cases}
\qquad
X(\omega)=\frac{2\sin(\omega T/2)}{\omega}.
\]
A finite rectangular pulse has a sinc-shaped spectrum with zeros at multiples of \(2\pi/T\).

\subsection{Properties}

\subsubsection{Fourier transform linearity}
\[
a x_1(t)+b x_2(t)\ \longleftrightarrow\ aX_1(\omega)+bX_2(\omega).
\]
Linearity allows spectra to be built from simpler transform pairs.

\subsubsection{Fourier transform time shift}
\[
x(t-t_0)\ \longleftrightarrow\ e^{-j\omega t_0}X(\omega).
\]
A time delay adds a linear phase term and does not change the magnitude spectrum.

\subsubsection{Fourier transform frequency shift}
\[
e^{j\omega_0t}x(t)\ \longleftrightarrow\ X(\omega-\omega_0).
\]
Multiplication by a complex sinusoid shifts the spectrum in frequency.

\subsubsection{Fourier transform differentiation}
\[
\frac{d^n x(t)}{dt^n}\ \longleftrightarrow\ (j\omega)^n X(\omega).
\]
Differentiation in time multiplies the spectrum by powers of \(j\omega\).

\subsubsection{Fourier transform convolution property}
\[
x(t)*h(t)\ \longleftrightarrow\ X(\omega)H(\omega).
\]
Convolution in time becomes multiplication in the Fourier domain.

\subsubsection{Fourier transform symmetry for real signals}
\[
x(t)\in\mathbb{R}\quad \Rightarrow\quad X(-\omega)=X^*(\omega).
\]
Real signals have conjugate-symmetric spectra. The magnitude is even and the phase is odd where the phase is defined.

\subsubsection{Fourier transform symmetry for even and odd signals}
\[
\begin{aligned}
x(t)\ \text{real and even} &\Rightarrow X(\omega)\ \text{real and even},\\
x(t)\ \text{real and odd} &\Rightarrow X(\omega)\ \text{imaginary and odd}.
\end{aligned}
\]
Signal symmetry strongly restricts the spectrum and is useful for true-false exam questions.
